% Industrial I - galpao leve com acessos e faixa de apoio refinados
\section{Projeto Industrial I — galpão de produção leve}
\newcommand{\GalpaoFinalBase}{%
\draw[arqExt] (0,0) rectangle (48,30);
\ArqParedeInt{0}{6}{38}{6}\ArqParedeInt{0}{8}{48}{8}
\foreach \x in {7,14,20,26,32,35}{\ArqParedeInt{\x}{0}{\x}{6}}
\foreach \x in {9,17,28,38}{\ArqParedeInt{\x}{8}{\x}{30}}
\ArqParedeInt{38}{0}{38}{8}
% portas da faixa de apoio
\ArqPortaHUp{1}{0}{1}
\ArqPortaHDown{3.8}{6}{0.9}\ArqPortaHDown{10.7}{6}{0.9}\ArqPortaHDown{16.0}{6}{0.9}\ArqPortaHDown{22.0}{6}{0.9}
\ArqPortaHDown{28.0}{6}{0.9}\ArqPortaHDown{33.0}{6}{0.9}\ArqPortaHDown{36.0}{6}{0.9}
% acessos do corredor aos setores
\ArqVaoH{2.0}{5.2}{8}\ArqVaoH{11.0}{14.2}{8}\ArqVaoH{20.5}{24.2}{8}\ArqVaoH{31.0}{34.2}{8}\ArqVaoH{41.0}{44.2}{8}
% passagens entre etapas
\foreach \x in {9,17,28,38}{\ArqVaoV{17}{21}{\x}}
% recebimento e docas
\ArqPortaCorrerV{13}{17}{0}\ArqPortaCorrerH{39}{42.2}{0}\ArqPortaCorrerH{43.5}{46.7}{0}
% apoio administrativo e tecnico
\ArqDesk{0.7}{1.2}{1.5}{0.6}\node[arqLabel] at (3.5,5.3){ADMIN.};
\ArqMesa{9}{2.5}{3}{1.1}\node[arqLabel] at (10.5,5.3){REFEITÓRIO / COPA};
\draw[arqEquip] (14.4,0.5) rectangle (19.6,5.5);\node[arqLabel] at (17,5.0){SALA ELÉTRICA};
\ArqMaquina{20.6}{1}{2}{1.2}{BANCADA}\node[arqLabel] at (23,5.3){MANUT.};
\ArqArmario{26.5}{0.6}{0.7}{4.7}\ArqArmario{30.7}{0.6}{0.7}{4.7}\node[arqLabel] at (29,5.3){ALMOXARIFADO};
\ArqArmario{32.35}{0.6}{0.55}{3.9}\ArqArmario{34.0}{0.6}{0.55}{3.9}\node[arqLabel] at (33.5,5.3){VESTIÁRIO};
\ArqVaso{36.1}{1.25}\ArqPia{35.35}{3.6}{0.75}{0.40}\ArqChuveiro{36.8}{3.4}{0.85}{1.3}\node[arqLabel] at (36.5,5.3){SANITÁRIOS};
% setores produtivos
\foreach \y in {10,14,18,22,26}{\draw[arqMob] (0.7,\y) rectangle (8.3,\y+1.3);}\node[arqLabel] at (4.5,29.2){RECEBIMENTO};
\foreach \y in {10,14,18,22,26}{\draw[arqMob] (9.7,\y) rectangle (16.3,\y+1.3);}\node[arqLabel] at (13,29.2){ESTOQUE MP};
\ArqMaquina{18.2}{11}{3.2}{2}{M1}\ArqMaquina{23}{11}{3.2}{2}{M2}\ArqMaquina{18.2}{18}{3.2}{2}{M3}\ArqMaquina{23}{18}{3.2}{2}{M4}\node[arqLabel] at (22.5,29.2){PROCESSO};
\foreach \x/\y in {29.2/11,33/11,29.2/18,33/18}{\ArqMaquina{\x}{\y}{2.6}{1.5}{POSTO}}\node[arqLabel] at (33,29.2){MONTAGEM};
\foreach \y in {10,14,18,22,26}{\draw[arqMob] (39,\y) rectangle (47.2,\y+1.3);}\node[arqLabel] at (43,29.2){ESTOQUE PA};
\node[arqLabel] at (43,7.2){EXPEDIÇÃO / DOCAS};\node[arqSub] at (24,7){CORREDOR LOGÍSTICO};
\draw[-{Stealth[length=3mm]},very thick,corNorma] (3,16)--(12,16)--(21,16)--(32,16)--(43,16);
\CotaH{0}{48}{30.22}{48,00 m}\CotaV{0}{30}{48.22}{30,00 m}}

\begin{frame}{Industrial I — acessos e apoio refinados}\centering\begin{tikzpicture}[scale=0.215,every node/.style={font=\tiny}]\GalpaoFinalBase\end{tikzpicture}\vspace{0.4mm}\scriptsize A faixa frontal agora tem funções definidas e portas próprias: administração, refeitório, elétrica, manutenção, almoxarifado, vestiário e sanitários. Vãos logísticos são separados dos acessos de pedestres.\end{frame}
\begin{frame}{Industrial I — lançamento elétrico sobre a planta}\centering\begin{tikzpicture}[scale=0.215,every node/.style={font=\tiny}]\GalpaoFinalBase\SimQD{16}{3}{QGBT}\SimQD{18.5}{3}{CCM}\foreach \x in {3,8,13,18,23,28,33,38,43}{\SimLuz{\x}{10.5}\SimLuz{\x}{23.5}}\end{tikzpicture}\end{frame}
